% Appendix of forgecert.tex.
\appendix
\section{Soundness}\label{app:proof}
Fix a box $B\subseteq D$ and the program of Eqs.~\eqref{eq:acc}--\eqref{eq:req}. For $x\in B$, let $r_n(x)$ be the value of neuron $n$ in $f$ and $\hat q_n(x)=t_n\,q_n(x)$ the value represented by the integer program (for inputs, $s_i q_i$).

\begin{definition}[Representation invariant]
A set of forms $\{(R_n,E_n)\}$ over symbols $\eta\in[-1,1]^K$ \emph{represents} a layer on $B$ if for every $x\in B$ there is $\eta$ with $R_n(\eta)=r_n(x)$ and $R_n(\eta)+E_n(\eta)=\hat q_n(x)$ for all neurons $n$ of the layer simultaneously.
\end{definition}
Fresh symbols introduced by a transfer function do not occur in earlier forms, so a witness $\eta$ for one layer extends to the next by choosing the fresh coordinates. Each lemma below shows that one transfer function preserves the invariant.

\begin{lemma}[Input]
$R_i=c_i+h_i\eta_i$ and $E_i=\rho_i\eta'_i$ with $\rho_i=s_i(\tfrac12+\delta_i)$ represent the input layer.
\end{lemma}
\begin{proof}
$\eta_i=(x_i-c_i)/h_i\in[-1,1]$ on $B$. The C code computes $t=\mathrm{fl}(x_i\cdot\mathrm{fl}(s_i^{-1}))$ and truncates $\mathrm{fl}(t\pm\tfrac12)$, so $|q_i-x_i/s_i|\le\tfrac12+\delta_i$, where $\delta_i$ bounds the two float32 roundings and the reciprocal (each at most $2^{-24}$ relative). The domain obligation $q^-\le q_i\le q^+$ rules out the clamp. Choose $\eta'_i=(s_iq_i-x_i)/\rho_i$.
\end{proof}

\begin{lemma}[Affine step]
If the inputs are represented, so are the forms $R'=WR+b$ and $E'=\Delta W R+\hat W E+\Delta b$ for the pre-activation values $r'$ and $\hat v$ (Eq.~\eqref{eq:eff}), with the same $\eta$.
\end{lemma}
\begin{proof}
Linear forms evaluate linearly: $R'(\eta)=\sum W r+b=r'(x)$ and $R'(\eta)+E'(\eta)=\sum\hat W\hat x+\hat b=\hat v(x)$. The second equality uses Eq.~\eqref{eq:eff}, which holds exactly because the domain obligations rule out overflow: by Lemma~\ref{lem:modular} for modular accumulators, by the partial-sum bound for classical ones, and by the product and $2^{62}$ bounds for the multiplies.
\end{proof}

\begin{lemma}[ReLU]
Let $Q=R+E$ and let $\widetilde{\mathrm{ReLU}}(Z)$ denote $\lambda Z+\mu+\mu\eta_{\mathrm{new}}$ with $\lambda=u/(u-l)$, $\mu=-\lambda l/2$ for an unstable $Z\in[l,u]$, the identity if $l\ge0$, and $0$ if $u\le0$. Both candidate error forms of \S\ref{sec:ai} preserve the invariant for $\mathrm{ReLU}(r')$ and $\mathrm{ReLU}(\hat v)$.
\end{lemma}
\begin{proof}
DeepZ property: for every $z\in[l,u]$, $\mathrm{ReLU}(z)=\widetilde{\mathrm{ReLU}}(z)$ for some value of the fresh symbol. (a) With independent fresh symbols for $R$ and $Q$, both relaxations are exact for suitable values at the same time. So $R''=\widetilde{\mathrm{ReLU}}(R')$ and $E''=\widetilde{\mathrm{ReLU}}(Q)-\widetilde{\mathrm{ReLU}}(R')$ give $R''+E''=\mathrm{ReLU}(\hat v)$. (b) $d=\mathrm{ReLU}(\hat v)-\mathrm{ReLU}(r')$ lies between $0$ and $\hat v-r'=E'(\eta)$ because ReLU is monotone and 1-Lipschitz. Each half-stable case is an identity: e.g.\ for $l\ge0$, $d=\max(E',-r')$. Hence $d$ lies in the computed interval $[a,b]$, and the fresh symbol takes the value $(2d-a-b)/(b-a)$. In the stable cases $d=0$ or $d=E'$ exactly.
\end{proof}

\begin{lemma}[Clamp]
For $Y$ the value after ReLU and $U=t\,q^+$, adding $-\widetilde{\mathrm{ReLU}}(Y-U)$ (and $\widetilde{\mathrm{ReLU}}(-U-Y)$ for signed outputs) to $E$ preserves the invariant for $\mathrm{clamp}(Y)$.
\end{lemma}
\begin{proof}
$\min(Y,U)=Y-\mathrm{ReLU}(Y-U)$ and $\max(Y,-U)=Y+\mathrm{ReLU}(-U-Y)$; apply the DeepZ property. Clamping commutes with rounding to multiples of $t$ because $U$ is such a multiple: $\mathrm{round}(\min(v,U))=\min(\mathrm{round}(v),U)$.
\end{proof}

\begin{lemma}[Rounding]
Adding a fresh symbol of radius $t/2$ preserves the invariant for $t\cdot\mathrm{rshift}(\mathrm{ReLU}(a)m,k)$. If the value is provably $0$ or $U$, no symbol is needed.
\end{lemma}
\begin{proof}
$\mathrm{rshift}(v,k)=\lfloor v/2^k+\tfrac12\rfloor$ and $t\cdot\mathrm{ReLU}(a)\,m\,2^{-k}=\mathrm{ReLU}(\hat v)$, so the represented output differs from $\mathrm{ReLU}(\hat v)$ by at most $t/2$. A value of $0$ or $U$ is a multiple of $t$ and rounds exactly.
\end{proof}

\begin{lemma}[Max-pooling]
Let the candidates of a window be represented by $(R_k,E_k)$, all with the same scale. The forms of \S\ref{sec:ai} represent $\max_k r_k$ and the integer program's $\max_k \hat q_k$.
\end{lemma}
\begin{proof}
The integer program compares codes of one scale $t>0$, so its result represents $\max_k \hat q_k$ exactly. A dropped candidate $k$ has $r_k<r_j$ and $\hat q_k<\hat q_{j'}$ for some $j,j'$ on all of $B$, so dropping it changes neither maximum. If $a$ dominates the remaining candidates in $R$, $\max_k r_k=r_a$; if $b$ dominates in $Q$, $\max_k\hat q_k=\hat q_b=R_b+E_b$; then $E'=Q_b-R_a$ is an identity. Otherwise a fresh symbol over $[\max_k\mathrm{lb},\max_k\mathrm{ub}]$ can take the value of the maximum, and $E'=Q'-R'$ holds by construction. For the alternative $E'$: if $\max r=r_a$ and $\max\hat q=\hat q_b$, then $e_a=\hat q_a-r_a\le \hat q_b-r_a\le \hat q_b-r_b=e_b$, so the difference of the maxima lies in $[\min_k e_k,\max_k e_k]$.
\end{proof}

\begin{lemma}[Graph]
If every tensor read by a layer is represented, so is its concatenation, and the affine step applies with taps over the concatenation. Batch normalisation folding, pooling lowering and residual additions are exact rewrites of the float model, and the checker performs them independently.
\end{lemma}

\begin{proof}[Proof of Theorem~\ref{thm:sound}]
By induction over the layers in topological order (each tensor's forms are kept until its last reader, and the buffer check guarantees the C program does the same with its data), the final forms represent the logits: for each $x\in B$ there is $\eta$ with $R_c(\eta)-R_j(\eta)=f_c(x)-f_j(x)$ and $Q_c(\eta)-Q_j(\eta)=t_L(y_c-y_j)$. The proof condition gives $f_c(x)-f_j(x)>0$, so $c=\arg\max f(x)$. It also gives $y_c-y_j>\theta_j/t_L$, i.e.\ $y_c\ge y_j+1$ for $j<c$ (integers) and $y_c>y_j$ for $j>c$, so the first-index argmax of the C code is $c$. The analysis evaluates forms in double precision with outward slack of $10^{-12}$ relative per bound. This is a documented assumption rather than directed rounding (\S\ref{sec:limits}).
\end{proof}
